% ============================================================
% Paper 6 - Positive Curvature of the Spectral Trace
%           at the Critical Line
%
% Author:   Ulrich Tehrani
% Date:     May 2026
% Version: v1.25
% Zenodo:   https://doi.org/10.5281/zenodo.19665790
% License:  CC BY 4.0
% ============================================================

\documentclass[11pt]{article}
\usepackage{cmap}
\usepackage[T1]{fontenc}
\usepackage[utf8]{inputenc}
\usepackage[a4paper,margin=1in]{geometry}
\usepackage{amsmath,amssymb,amsthm,mathtools}
\usepackage{xcolor}
\usepackage{hyperref}
\usepackage{enumitem}
\usepackage{setspace}
\usepackage{booktabs}
\usepackage{array}
\usepackage{graphicx}

\hypersetup{colorlinks=true,
  linkcolor=blue!50!black,
  urlcolor=blue!50!black,
  citecolor=blue!50!black}
\setlength{\parskip}{0.5em}
\setlength{\parindent}{0pt}
\onehalfspacing
\emergencystretch=1em

% --------------------------------------------------------
% Theorem environments (numbered within section)
% --------------------------------------------------------
\newtheorem{theorem}{Theorem}[section]
\newtheorem{proposition}[theorem]{Proposition}
\newtheorem{lemma}[theorem]{Lemma}
\newtheorem{corollary}[theorem]{Corollary}
\newtheorem{definition}[theorem]{Definition}
\newtheorem{assumption}[theorem]{Assumption}
\newtheorem{observation}[theorem]{Observation}
\newtheorem{conjecture}[theorem]{Conjecture}
\newtheorem*{remark}{Remark}
\newtheorem{openproblem}[theorem]{Open Problem}

% Heuristic environment (clearly marked, never used in proofs)
\newenvironment{heuristic}[1]{%
  \par\smallskip\noindent\textit{[Heuristic: #1]}\quad}{%
  \par\smallskip}

% --------------------------------------------------------
% Convention box — clean ruled box, no mdframed
% --------------------------------------------------------
\newenvironment{conventionbox}{%
  \medskip
  \noindent\rule{\linewidth}{0.4pt}\smallskip\par\noindent
}{%
  \par\smallskip\noindent\rule{\linewidth}{0.4pt}
  \medskip
}

% --------------------------------------------------------
% Series-wide macros
% --------------------------------------------------------
\newcommand{\Hstr}{H_{\mathrm{str}}}
\newcommand{\Hnull}{H_{\mathrm{null}}}
\newcommand{\Estr}{E_{\mathrm{str}}}
\newcommand{\Espec}{E_{\mathrm{spec}}}
\newcommand{\etaorig}{\eta_{\mathrm{orig}}}
\newcommand{\Tren}{T_{\mathrm{ren}}}
\newcommand{\Tt}{\widetilde{T}}
\newcommand{\Gun}{G^{\mathrm{un}}}
\newcommand{\Hlocal}{H_{\mathrm{local}}}

% Smoothed zero sums and trace objects
\newcommand{\DSEL}{D_{\mathrm{SEL}}}
\newcommand{\Ediag}{E_{\mathrm{diag}}}
\newcommand{\Ddiag}{D_{\mathrm{diag}}}

% Constants (always use subscript — never plain C)
\newcommand{\Ceta}{C_{\eta}}
\newcommand{\CT}{C_{T}}

% --------------------------------------------------------
% Title
% --------------------------------------------------------
\title{\vspace{-1.5em}\bfseries Positive Curvature of the Spectral Trace
at the Critical Line}
\author{Ulrich Tehrani}
\date{Research Note \textperiodcentered\ August 2026
      \textperiodcentered\ v1.25}

% ============================================================
\begin{document}
\maketitle
% ============================================================

\begin{abstract}
Building on the exact trace formula $\mathrm{Tr}(\Tt(\sigma))=\DSEL-O(\sigma)$
established in~\cite{Paper5}, we give a structural analysis of the
second-order curvature of the spectral trace at the critical line
and reduce its sign to a single weighted prime--zero sum~$B$.
The central quantity is the direct curvature sum
$B=\sum_{k,p} e^{-\varepsilon^2\gamma_k^2}(\gamma_k\log p)^2\cos(\gamma_k\log p)$,
which satisfies the algebraic identity $O''(\tfrac12)=-2B$, together with a
Guinand--Weil decomposition of a smoothed prime--zero sum $\widetilde Z_p(\varepsilon)$
whose leading term is proved strictly negative.
No hypothetical input is used: no Riemann Hypothesis, no Hilbert--P\'olya
postulate, no GUE conjecture.

\textbf{What is proved.}
For every prime $p$ and every $\varepsilon>0$, the leading term of the
Guinand--Weil decomposition satisfies
$\mathrm{Main}_p(\varepsilon)<0$ (Theorem~\ref{thm:main-neg}), while the
other-prime error is non-positive (Theorem~\ref{thm:other}).
The algebraic identity $O''(\tfrac12)=-2B$ connecting curvature to bias
is derived from the trace formula of~\cite{Paper5}
(Proposition~\ref{prop:curv-bias}).
The curvature sign of the spectral trace at $\sigma=\tfrac12$
thus reduces to the sign of a single scalar quantity~$B$.
The decomposition
$\widetilde Z_p(\varepsilon)=\mathrm{Main}_p+\Gamma_p+\mathrm{Const}_p
+\mathrm{Err}_{p^k}+\mathrm{Err}_{\mathrm{other}}$
is the classical Guinand--Weil explicit formula applied to an even,
Gaussian-damped test function; its derivation is recalled in~\S\,\ref{sec:gw}.

\textbf{What is numerical.}
The gamma factor contribution is subleading, with
$|\Gamma_p(\varepsilon)|/|\mathrm{Main}_p(\varepsilon)|$ numerically
consistent with linear scaling in~$\varepsilon$
(Theorem~\ref{thm:gamma}).
The zero-sum truncation tail at $(\varepsilon,N)=(0.05,100)$ is
verified numerically to be negligible, with ratio bounded by
$3\times10^{-61}$ relative to the main term (Theorem~\ref{thm:trunc}).
The proxy ratio $r^+_{p,N}(\varepsilon)$ appears unimodal on the
tested $\varepsilon$-grid:
it exceeds~$1$ for $p\in\{37,53\}$ near $\varepsilon=0.05$ and
is consistent with decrease toward $r_p^{\infty}\le\tfrac12$
as $\varepsilon\to 0$
(Theorem~\ref{thm:ratio-signcross}(a), Remark~\ref{rem:rpinf}).
On the sampled grid, the observed transition occurs between the
tested values $\varepsilon=0.020$ and $\varepsilon=0.025$, with
$Z_{p,N}^+(\varepsilon)<0$ at the tested $\varepsilon=0.020$
(ordinate proxy $Z_{p,N}^+=\Re Z_{p,\varepsilon,N}$)
for every prime $p\le 53$ (Theorem~\ref{thm:ratio-signcross}(b)).
The integrated bias is negative at reference parameters,
$B_{\mathrm{int}}^+(0.05,100)=-42.21<0$ at
$(\kappa,\varepsilon,N)=(53,0.05,100)$ (Theorem~\ref{thm:intbias}).
At the same reference parameters the direct curvature sum equals
$B=-19\,342.5<0$, which is equivalent to
$O''(\tfrac12)=+38\,685>0$ by the curvature--bias identity
(Corollary~\ref{cor:curvature}).

\textbf{What is conditional.}
If, in addition to the numerically established inequality $B<0$ at the
reference parameters, the stationarity condition $O'(\tfrac12)=0$
holds at the same parameters, then $\sigma=\tfrac12$ is a strict local
minimum of $O$ and $\mathrm{Tr}(\Tt(\sigma))$ attains a strict local
maximum at $\sigma=\tfrac12$ at the reference parameters
(Remark following Corollary~\ref{cor:curvature}).
The conditional statement is clearly marked;
a weaker first-derivative cancellation estimate is formulated in
Open Problem~\ref{op:stat}; exact stationarity $O'(\tfrac12)=0$
remains open.

\textbf{What remains open.}
Five questions separate the present work from a full curvature proof,
each precisely named in~\S\,\ref{sec:open}.
Problem~\ref{op:stat} (\emph{First-derivative cancellation}) asks for
a nontrivial cancellation estimate for the prime exponential sums
$S_\kappa(\gamma_k)$ strong enough to give a genuine logarithmic
saving in $|O'(\tfrac12)|$ relative to the trivial triangle bound.
Problem~\ref{op:bias-asymp} (\emph{Asymptotic pointwise bias}) asks
for an analytic proof that $Z_p^\infty(\varepsilon)<0$ for every prime
$p\le\kappa$ and sufficiently small $\varepsilon>0$, together with a
finite-$N$ transfer to $Z_{p,N}^+(\varepsilon)$.
Problem~\ref{op:uniformity} (\emph{Uniformity of positive curvature})
asks whether the inequalities $B<0$ and $O''(\tfrac12)>0$ persist at
parameter points $(\kappa,\varepsilon,N)$ other than the reference
values tested here.
Problem~\ref{op:transfer} (\emph{Integrated-to-direct transfer}) asks
whether pointwise or integrated bias implies the negativity of $B$.
Problem~\ref{op:weil-bridge} (\emph{Lagarias--Weil bridge}) asks whether
the curvature statement transfers to the Weil-positivity framework
of~\cite{Lagarias}.

\textbf{Structure.}
The formal core in \S\S\,\ref{sec:gw}--\ref{sec:curvcor} is non-circular;
analytic identities, numerical observations, and conditional statements
are explicitly separated.
No use of the Riemann Hypothesis, the GUE conjecture, the Montgomery
pair correlation conjecture, or the Hilbert--P\'olya postulate is made
anywhere in \S\S\,\ref{sec:gw}--\ref{sec:curvcor}.
\end{abstract}

\vspace{2em}

% --------------------------------------------------------
% Notation and Conventions
% --------------------------------------------------------
\section*{Notation and Conventions}
\label{sec:notation}

\emph{Critical line.}
Throughout, the critical line means $\Re(s)=\tfrac12$.
The trace parameter $\sigma\in\mathbb{R}$ is evaluated at
$\sigma=\tfrac12$ unless stated otherwise.

\medskip
\emph{Global parameters.}
The reference values are $\kappa=53$ (prime cutoff, yielding the
$16$ active primes $p\le\kappa$), $\varepsilon=0.05$ (Gaussian
regularisation parameter), and $N=100$ (number of Riemann ordinates
$\gamma_k$ employed).
These values are fixed throughout
\S\S\,\ref{sec:gw}--\ref{sec:curvcor} unless explicitly varied.
The associated spectral weight sums are
\[
W_{\varepsilon,N}:=\sum_{k=1}^N e^{-\varepsilon^2\gamma_k^2},
\qquad
P(\kappa):=\sum_{p\le\kappa}\log p.
\]

\medskip
\emph{Main objects.}
The finite truncation
\[
Z_p \;=\; Z_{p,\varepsilon,N} \;:=\;
\sum_{k=1}^{N} e^{-\varepsilon^2\gamma_k^2}\,p^{i\gamma_k}
\]
is the numerical object of~\cite{Paper4,Paper5}.
The $\gamma_k^2$-weighted variant
\[
Z_p^{(2)} \;:=\; \sum_{k=1}^{N} \gamma_k^2\,
e^{-\varepsilon^2\gamma_k^2}\,p^{i\gamma_k}
\]
carries the curvature weighting needed for the bias
decomposition; the dual form of the curvature sum is
$B = \sum_p (\log p)^2\,\mathrm{Re}\,Z_p^{(2)}$.
The associated Guinand--Weil object
\[
\widetilde Z_p(\varepsilon) \;:=\; \sum_{\rho}
h^{\mathrm{even}}_{p,\varepsilon}\!\left(\tfrac{\rho-1/2}{i}\right),
\qquad
h^{\mathrm{even}}_{p,\varepsilon}(u)=e^{-\varepsilon^2 u^2}\cos(u\log p),
\]
sums over all non-trivial zeros of $\zeta$.
$\widetilde Z_p(\varepsilon)$ is real-valued by the quartet symmetry
$\{\rho,\bar\rho,1-\rho,1-\bar\rho\}$ of the non-trivial zeros.
The finite positive-ordinate cosine proxies
$Z_{p,N}^+(\varepsilon)=\Re Z_{p,\varepsilon,N}$ and $\Re Z_p^{(2)}$
are real by construction (cosine kernel).
The phasor sums $Z_p$ and $Z_p^{(2)}$ themselves are complex;
only their real parts enter the cosine-kernel computations used here.
The identity $\widetilde Z_p^{\mathrm{full}} = 2\,Z_p^{+}$ holds when
all non-trivial zeros lie on the critical line; in general
$\widetilde Z_p^{\mathrm{full}} = 2\,Z_p^{+} + E_p^{\mathrm{off}}$
where $E_p^{\mathrm{off}}$ captures off-critical-line contributions
(see Paper~5, Open Problem~6.5 (GW bridge)).
All numerical computations in this paper use the ordinate proxy
$Z_p^{+}(\varepsilon) = \Re Z_{p,\varepsilon,N}$ directly.
The Guinand--Weil object $\widetilde Z_p$ motivates the decomposition
but the actual numerical values refer to $Z_p^{+}$.
At $(\varepsilon,N)=(0.05,100)$ the truncation tail $R_{p,N}$
is negligible by Theorem~\ref{thm:trunc}.
The direct curvature sum
\[
B \;:=\; \sum_{k,p} e^{-\varepsilon^2\gamma_k^2}\,
(\gamma_k\log p)^2\,\cos(\gamma_k\log p)
\]
enters the algebraic identity $O''(\tfrac12)=-2B$
(Proposition~\ref{prop:curv-bias}); the integrated bias is
$B_{\mathrm{int}}^+(\varepsilon,N)\;:=\;\sum_p (\log p)^2\,Z_{p,N}^+(\varepsilon)$.

\medskip
\emph{Analytical foundations.}
The identity $O''(\tfrac12)=-2B$ (Proposition~\ref{prop:curv-bias})
is obtained by direct differentiation of the trace oscillatory sum
of Paper~5.
The Guinand--Weil decomposition in \S\,\ref{sec:gw} uses the classical
explicit formula framework of Weil~\cite{Weil} and Guinand~\cite{Guinand}.
The truncation estimate in Theorem~\ref{thm:trunc} uses
Riemann--von Mangoldt zero-counting
estimates~\cite[Thm.\,5.8]{IwaniecKowalski04} and Gaussian tail bounds.

\medskip
\emph{Numerical computations.}
All numerical results were computed with \texttt{mpmath}
at 50~decimal digits of working precision;
Riemann zeros $\gamma_k$ were evaluated via
\texttt{mpmath.zetazero}.
Verification code is available at~\cite{Code}
(\url{https://github.com/utehrani/analysislab-nt}).

\medskip
\emph{Not to be confused with.}
$B$ (curvature sum, weight $(\gamma_k\log p)^2$) and $B_{\mathrm{int}}$
(integrated bias, weight $(\log p)^2$) differ in their weighting and
in the object summed; both are negative at the reference parameters,
but neither implies the other (Open Problem~\ref{op:transfer}).
$Z_p$ (finite truncation, level~$N$, sum over positive ordinates
$k=1,\ldots,N$, in contract notation $Z_p^{+}$),
$Z_p^{(2)}$ (with $\gamma_k^2$-weighting, used in~$B$), and
$\widetilde Z_p$ (Guinand--Weil object, sum over all non-trivial zeros,
in contract notation $\widetilde Z_p^{\mathrm{full}}$) are three
formally distinct objects. For $h^{\mathrm{even}}$ and $N\to\infty$, the identity
$\widetilde Z_p^{\mathrm{full}}(\varepsilon) = 2\,Z_p^{+}(\varepsilon)$
holds conditionally on all non-trivial zeros being on the critical line;
without this assumption, an off-critical contribution $E_p^{\mathrm{off}}$
must be added (cf.\ Paper~5, Open Problem~6.5 (GW bridge)).
Numerical computations therefore refer to the ordinate proxy
$Z_p^{+}(\varepsilon)$, not the full GW object~$\widetilde Z_p^{\mathrm{full}}$.
$Z_p^{(2)}$ carries the additional spectral weight absent from~$Z_p$.
The diagonal-energy constants $D\approx 9.470$ of~\cite{Paper2} and
$\DSEL\approx 10.985$ of~\cite{Paper5} share a name across the series
but are defined in different operator contexts and must not be
identified.

\newpage

% -------------------------------------------------------
\section{Introduction}
\label{sec:intro}

\subsection{Background and Motivation}

This paper is the sixth in a series developing an operator-theoretic
framework for studying the distribution of the non-trivial zeros of
the Riemann zeta function $\zeta(s)$.
The series is built around a finite-dimensional Hilbert-space model in
which prime numbers mediate coupling between zeros, encoded in an
explicit operator family.

\textbf{Paper~1}~\cite{Paper1} introduced the local curvature function
$\Hlocal(\sigma,\kappa)=\psi_1(\sigma)+\sum_{p\le\kappa}V_p(\sigma)$, identifying
$\sigma=\tfrac12$ as a phase boundary through the divergence
$\Hlocal(\tfrac12,\kappa)\sim 2(\log\kappa)^2$.

\textbf{Paper~2}~\cite{Paper2} connected this local curvature to the
Weil explicit functional, constructing admissible test functions
for which the curvature divergence becomes a structural feature of
the Weil energy.

\textbf{Paper~3}~\cite{Paper3} made the geometric structure explicit:
two finite-dimensional Hilbert spaces $\Hstr$ (over primes) and
$\Hnull$ (over zeros) connected by a linear coupling map $\Phi$, with
loop operator $T=\Phi^*\circ\Phi$ and a first numerical verification of the
energy asymmetry $\etaorig>0$.

\textbf{Paper~4}~\cite{Paper4} introduced the dual operator
$\Tt=\Phi\circ\Phi^*$ on $\Hnull$ --- the Tehrani operator --- and
established a conditional proof of energy asymmetry under the
remainder-control assumption $(E_{\mathrm{rem}})$, motivated by
prime-phase cancellation estimates.

\textbf{Paper~5}~\cite{Paper5} extended this to a one-parameter
family $\Tt(\sigma)=\Phi(\sigma)\circ\Phi(\sigma)^*$ and proved the
exact spectral trace identity
\[
\mathrm{Tr}(\Tt(\sigma)) \;=\; \DSEL - O(\sigma),
\qquad
O(\sigma) \;=\; \tfrac12\sum_{k,p} e^{-\varepsilon^2\gamma_k^2}
                \cos(2\sigma\gamma_k\log p),
\]
where $\DSEL$ is a $\sigma$-independent constant.
Four independent numerical signatures concentrating at $\sigma=\tfrac12$
were documented in~\cite{Paper4,Paper5}, and the curvature-bias quantity
$B=\sum_p(\log p)^2\,\mathrm{Re}(Z_p^{(2)})$ was introduced with the reference
value $B=-19\,342.5$ at $(\kappa,\varepsilon,N)=(53,0.05,100)$.

\textbf{The present paper} addresses the analytic question left open
by~\cite{Paper5}: \emph{is there a structural reason, independent of
the Riemann Hypothesis, why $\sigma=\tfrac12$ is distinguished by the
trace?}
We reduce this question at second order to the sign of the direct
curvature sum $B$ via the algebraic identity $O''(\tfrac12)=-2B$,
and verify the positive-curvature sign numerically at the reference
parameters.
At the reference parameters $B$ is numerically negative, hence
$O''(\tfrac12)$ is numerically positive.
The paper establishes this curvature statement together with a
Guinand--Weil decomposition of a related smoothed prime--zero sum,
recalls four further numerical signatures from Papers~4--5 concentrating
at $\sigma=\tfrac12$,
and formulates five open problems constituting the completion path
toward a theorem-level selection principle.

The operator $\Tt(\sigma)$ should be viewed neither as a
Hilbert--P\'olya spectral realisation in the sense of
Connes~\cite{Connes} or Berry--Keating~\cite{BerryKeating}, nor as a
de~Branges structure-function criterion~\cite{deBranges}; it is a
finite-cutoff operator packaging Gaussian-smoothed explicit-formula
data, closest in spirit to Hilbert-space reformulations of the
explicit formula such as those of Burnol~\cite{Burnol}.

The heuristic labels ``prime field'' and ``field strength'' used in
the series denote a finite-cutoff coupling picture in which primes
mediate pairwise interactions between zero ordinates via explicit
Gaussian-smoothed phasors; this should be distinguished from the
periodic-orbit dictionary of Berry--Keating~\cite{BerryKeating},
the adelic trace-formula picture of Connes~\cite{Connes}, and the
Frobenius/cohomology picture of function-field zeta functions.

\subsection{Main Results}

The paper contains three main components.

The first is the curvature statement at the critical line.
At the reference parameters $(\kappa,\varepsilon,N)=(53,0.05,100)$,
direct numerical evaluation gives $B=-19\,342.5<0$, which is equivalent
via the algebraic identity of Proposition~\ref{prop:curv-bias} to
$O''(\tfrac12)=-2B=+38\,685>0$
(Corollary~\ref{cor:curvature}).
A strict local minimum of $O$ at $\sigma=\tfrac12$ would require
additionally that $O'(\tfrac12)=0$. At the reference parameters,
$O'(\tfrac12)=+2.475\ne0$, so $\sigma=\tfrac12$ is not an exact
local minimum of $O$ at the reference finite model.
A sufficiently small stationarity bound, together with a local
convexity lower bound for $O''$ near $\sigma=\tfrac12$, would
locate a local minimum near $\sigma=\tfrac12$
(Problem~\ref{op:stat}).

The curvature identity $O''(\tfrac12)=-2B$ is structurally analogous
to the Li/Bombieri--Lagarias sign
criteria~\cite{Li,BombieriLagarias} in that it turns
explicit-formula data into a sign condition, but it is not a
derivative of $\xi$ at $s=1$ and not a Weil-positivity theorem; it
is a local second-order statement in the auxiliary trace parameter
$\sigma$ at the symmetry point~$\tfrac12$.

Bombieri~\cite{Bombieri} studies the positivity of the
Weil quadratic functional $Q(g) = W(g\ast\tilde g)$ truncated
to $L^2[-t,t]$, proving positivity for small $t$ and showing that
off-critical zeros contribute exactly half a negative eigenvalue
each.
The present paper does not truncate the functional in a spatial
variable: our operator $\Tt(\sigma)=\Phi(\sigma)\Phi(\sigma)^*$
is a Gram-type finite-rank matrix indexed by explicit
Gaussian-damped prime phasors at spectral cutoff $\kappa$ and
regularisation~$\varepsilon$, and our central object is the
\emph{second derivative} $O''(\tfrac12) = -2B$ of an exact trace
(not a spectral positivity condition on $Q$).
The Bombieri quadratic form is the appropriate ambient framework
for a possible future construction of the Weil transfer
operator~$W_{\kappa,\varepsilon}$.

The second is the Guinand--Weil bias analysis of~\S\S\,\ref{sec:gw}--\ref{sec:const}.
The smoothed prime--zero sum $\widetilde Z_p(\varepsilon)$ is decomposed into
five contributions, of which the leading one is proved negative
(Theorem~\ref{thm:main-neg}), and the other-prime contribution is proved non-positive
(Theorem~\ref{thm:other}), while the zero-ordinate truncation
tail is numerically verified, with analytic tail control, to be
negligible at the reference parameters
(Theorem~\ref{thm:trunc}).
The gamma contribution is shown numerically to be subleading of order
$\varepsilon$ (Theorem~\ref{thm:gamma}).
The proxy ratio $r^+_{p,N}(\varepsilon)$ is unimodal on the tested grid
(Theorem~\ref{thm:ratio-signcross}(a), Remark~\ref{rem:rpinf}), and a sign-crossover
is localised in $(0.020,\,0.025)$
(Theorem~\ref{thm:ratio-signcross}(b)).
The integrated bias
$B_{\mathrm{int}}^+(\varepsilon,N)=\sum_{p\le\kappa}(\log p)^2\,
Z_{p,N}^+(\varepsilon)$ is numerically negative at the
reference parameters, with $B_{\mathrm{int}}^+(0.05,100)=-42.21$
(Theorem~\ref{thm:intbias}).

The third is the structural separation made explicit in~\S\,\ref{sec:intbias}.
The quantities $B$, $B_{\mathrm{int}}(\varepsilon)$, and the pointwise
sign of $Z_{p,N}^+(\varepsilon)$ are structurally distinct
witnesses of the same qualitative phenomenon: $B$ carries the
$(\gamma_k\log p)^2$-weighting, $B_{\mathrm{int}}$ only the
$(\log p)^2$-weighting, and the pointwise statement carries no
weighting at all.
$B$ and $B^+_{\mathrm{int}}$ are numerically negative at the reference parameters,
and the pointwise bias $Z_{p,N}^+(0.05)<0$ holds for $14$ of the $16$ primes; the
analytical relationship between them is left open as
Problem~\ref{op:transfer}.
The five open problems of~\S\,\ref{sec:open} list the steps remaining,
in order of what we consider tractability:
Problem~\ref{op:stat} (first-derivative cancellation) is the most tractable and,
together with the numerically established $B<0$ and local convexity,
would locate a strict local minimum of $O$ near $\sigma=\tfrac12$.
An exact strict local minimum at $\sigma=\tfrac12$ requires the
stronger condition $O'(\tfrac12)=0$.

\subsection{Reader Contract}

This paper \emph{proves}, by direct computation and classical estimates
of the Guinand--Weil type, that the main term is negative
(Theorem~\ref{thm:main-neg}), that the other-prime error is
non-positive (Theorem~\ref{thm:other}), and the algebraic
curvature--bias identity $O''(\tfrac12)=-2B$
(Proposition~\ref{prop:curv-bias}).
These are unconditional mathematical statements.

This paper \emph{establishes numerically}, at the reference parameters
$(\kappa,\varepsilon,N)=(53,0.05,100)$,
that the truncation tail is negligible at the reference parameters
(Theorem~\ref{thm:trunc}, numerically verified at
$(\varepsilon,N)=(0.05,100)$),
that the gamma contribution is
subleading of order $\varepsilon$ (Theorem~\ref{thm:gamma}), that
the proxy ratio $r^+_{p,N}(\varepsilon)$ is unimodal on the
tested grid and is consistent with decrease toward $r_p^{\infty}\le\tfrac12$
(Theorem~\ref{thm:ratio-signcross}(a),
Remark~\ref{rem:rpinf}), and that the pointwise
bias $Z_{p,N}^+(\varepsilon)<0$ holds on the tested subgrid at
$\varepsilon=0.020$ (via ordinate proxy; Theorem~\ref{thm:ratio-signcross}), that
$B_{\mathrm{int}}^+(0.05,100)=-42.21<0$ (Theorem~\ref{thm:intbias}), and that
the positive-curvature statement $O''(\tfrac12)=+38\,685>0$ holds at
the reference parameters (Corollary~\ref{cor:curvature}).
These statements are verified on a finite grid and are clearly labelled
[numerical] in their theorem headings.

This paper \emph{states conditionally} that exact stationarity
$O'(\tfrac12)=0$, together with the numerically established
$O''(\tfrac12)>0$ at the reference parameters, would make
$\sigma=\tfrac12$ a strict local minimum of $O$
(Remark following Corollary~\ref{cor:curvature}).
A stationarity bound $|O'(\tfrac12)|\ll W_{\varepsilon,N}\cdot P(\kappa)$ would
become a localisation statement only when additionally combined with
a local convexity lower bound and the smallness condition stated in
Problem~\ref{op:stat}.
No part of the unconditional claims above depends on either condition.

This paper \emph{leaves open} five well-posed questions named in
\S\,\ref{sec:open}: Problem~\ref{op:stat} (First-derivative cancellation),
Problem~\ref{op:bias-asymp} (Asymptotic pointwise bias),
Problem~\ref{op:uniformity} (Uniformity of positive curvature),
Problem~\ref{op:transfer} (Integrated-to-direct transfer), and
Problem~\ref{op:weil-bridge} (Lagarias--Weil bridge).
The first four are formulated entirely within classical analytic number
theory; the fifth names the remaining bridge to Weil positivity in
the formulation of~\cite{Lagarias}.

% -------------------------------------------------------
\section{The Guinand--Weil Decomposition}
\label{sec:gw}

Throughout, $p$ denotes a rational prime, $\varepsilon>0$ a smoothing
parameter, and $\rho=\beta+i\gamma$, $0<\beta<1$, a non-trivial zero
of the Riemann zeta function.
For critical-line zeros, $(\rho-\tfrac12)/i=\gamma$; off the critical
line, the argument is complex. The real-valuedness of $\widetilde Z_p$
follows from the quartet symmetry
$\{\rho,\bar\rho,1-\rho,1-\bar\rho\}$ of the non-trivial zeros.
We fix a cutoff $\kappa\ge 2$ and consider only primes $p\le\kappa$.
For numerical reference we use $(\kappa,\varepsilon,N)=(53,0.05,100)$,
where $N$ is the number of zeros used in the finite approximation;
the first $16$ primes lie below $\kappa=53$.
The Guinand--Weil explicit formula
(see e.g.\ \cite[\S\,5.5]{IwaniecKowalski04}) is the key analytic tool;
its admissibility conditions on the test function are recalled below.

\begin{definition}[Smoothed zero sum]
\label{def:smoothed}
Let
\[
Z_p \;=\; Z_{p,\varepsilon,N} \;:=\; \sum_{k=1}^{N} e^{-\varepsilon^2\gamma_k^2}\,
p^{i\gamma_k},
\qquad
\mathrm{Re}\,Z_p \;=\; \sum_{k=1}^{N} e^{-\varepsilon^2\gamma_k^2}
                        \cos(\gamma_k\log p).
\]
This is the finite numerical object studied in~\cite{Paper4,Paper5}.
The associated object obtained from the Guinand--Weil explicit formula is
\[
\widetilde Z_p(\varepsilon) \;:=\; \sum_{\rho} h^{\mathrm{even}}_{p,\varepsilon}
\!\left(\tfrac{\rho-\tfrac12}{i}\right),
\qquad
h^{\mathrm{even}}_{p,\varepsilon}(u) \;=\; e^{-\varepsilon^2 u^2}
                                            \cos(u\log p).
\]
\end{definition}

The Gaussian $h^{\mathrm{even}}_{p,\varepsilon}$ is admissible
in the sense of Guinand--Weil (entire of order~2, Schwartz-class
decay, even, real-valued; cf.\
Carneiro--Littmann--Vaaler~\cite{CarneiroLittmannVaaler13}
for the Beurling--Selberg Gaussian-subordination framework;
and Carneiro--Milinovich--Soundararajan~\cite{CarneiroMilinovichSoundararajan19}
for Fourier-optimization methods in the explicit-formula context).
Since $h^{\mathrm{even}}_{p,\varepsilon}$ is even and decays rapidly,
$\widetilde Z_p(\varepsilon)$ is real-valued and absolutely convergent.
Its Fourier transform is
\[
\widehat{h}_{p,\varepsilon}(x) \;=\; \frac{\sqrt{\pi}}{2\varepsilon}
\!\left[\exp\!\left(-\tfrac{(\log p-2\pi x)^2}{4\varepsilon^2}\right)
      +\exp\!\left(-\tfrac{(\log p+2\pi x)^2}{4\varepsilon^2}\right)
\right],
\]
which is strictly positive on $\mathbb R$ and concentrates at
$\pm\log p/(2\pi)$ as $\varepsilon\to 0$.
We use the Fourier convention
$\widehat{h}(x)=\int_{-\infty}^{\infty}h(u)\,e^{-2\pi iux}\,du$,
matching the arguments $\log n/(2\pi)$ in the Guinand--Weil formula.

\begin{proposition}[Guinand--Weil decomposition convention; classical formula recalled]
\label{prop:gw}
For every prime $p$ and every $\varepsilon>0$,
\[
\widetilde Z_p(\varepsilon) \;=\; \mathrm{Main}_p(\varepsilon)
                       + \Gamma_p(\varepsilon)
                       + \mathrm{Const}_p(\varepsilon)
                       + \mathrm{Err}_{p^k}(\varepsilon)
                       + \mathrm{Err}_{\mathrm{other}}(\varepsilon),
\]
where the five contributions are defined by the standard partition of
the prime sum and the archimedean terms of the explicit formula.
Explicitly,
\[
\mathrm{Main}_p(\varepsilon) \;=\; -\frac{1}{2\pi}\cdot\frac{\log p}{\sqrt p}
\!\left[\widehat{h}_{p,\varepsilon}\!\left(\tfrac{\log p}{2\pi}\right)
      +\widehat{h}_{p,\varepsilon}\!\left(-\tfrac{\log p}{2\pi}\right)
\right],
\]
\[
\Gamma_p(\varepsilon) \;=\; \frac{1}{2\pi}\int_{-\infty}^{\infty}
e^{-\varepsilon^2 t^2}\cos(t\log p)\,
\mathrm{Re}\,\psi\!\left(\tfrac14 + \tfrac{it}{2}\right)\,dt,
\]
\[
\mathrm{Const}_p(\varepsilon) \;=\; \widetilde Z_p(\varepsilon)
\,-\, \mathrm{Main}_p(\varepsilon)
\,-\, \Gamma_p(\varepsilon)
\,-\, \mathrm{Err}_{p^k}(\varepsilon)
\,-\, \mathrm{Err}_{\mathrm{other}}(\varepsilon),
\]
where $\psi$ denotes the digamma function.
The prime-power error $\mathrm{Err}_{p^k}$ collects the contribution
from $n=p^m$, $m\ge 2$ (higher powers of the fixed prime $p$).
The remaining-prime error $\mathrm{Err}_{\mathrm{other}}$ collects
the contribution from $n=q^m$ with $q\ne p$ prime and $m\ge 1$
(all prime powers not involving $p$).
Both are standard sums
$\sum_n\Lambda(n)\widehat{h}_{p,\varepsilon}(\log n/(2\pi))/\sqrt n$.
\end{proposition}

This proposition records the decomposition convention induced by the
classical Guinand--Weil explicit formula. The present paper does not
supply a new proof of the explicit formula; it invokes the standard
formula with the stated Fourier normalisation and defines
$\mathrm{Const}_p$ as the residual after removing the other displayed
contributions.

\begin{remark}
At the test-function evaluations corresponding to the zeta
singularity at $s=1$ and its functional-equation reflection at
$s=0$, the completed explicit formula contributes the value
$h^{\mathrm{even}}_{p,\varepsilon}(u)$ evaluated at $u=\pm i/2$.
A direct computation gives
\[
h^{\mathrm{even}}_{p,\varepsilon}\!\left(\tfrac{i}{2}\right)
+ h^{\mathrm{even}}_{p,\varepsilon}\!\left(-\tfrac{i}{2}\right)
\;=\; 2\,e^{\varepsilon^2/4}\cosh\!\left(\tfrac{\log p}{2}\right)
\;=\; e^{\varepsilon^2/4}\!\left(\sqrt p + \tfrac{1}{\sqrt p}\right).
\]
This is the hyperbolic cosine, not the circular cosine: the formal
substitution $\cos(i\log p/2)=\cosh(\log p/2)$ is essential.
Replacing $\cosh$ by $\cos$ would yield a quantity that is negative
for $p\ge 29$, which is analytically impossible since the left-hand
side equals $\sqrt p + 1/\sqrt p>2$ for every $p\ge 2$.
The correct $\cosh$ form is used throughout what follows.
See Edwards~\cite[\S1.8, \S1.16, and \S3.6]{Edwards} for the classical derivation
of the pole contributions at $s=0,1$.
\end{remark}

The remainder of the paper studies the five terms of the decomposition
in turn.

% -------------------------------------------------------
\section{The Main Term}
\label{sec:main}

\begin{theorem}[Main term negativity; proved]
\label{thm:main-neg}
For every prime $p$ and every $\varepsilon>0$,
\[
\mathrm{Main}_p(\varepsilon) \;=\; -\frac{\log p}{2\sqrt\pi\,\varepsilon\,\sqrt p}
\cdot\bigl(1 + e^{-(\log p)^2/\varepsilon^2}\bigr)
\;<\; 0.
\]
\end{theorem}

\begin{proof}
The Fourier transform $\widehat{h}_{p,\varepsilon}(x)$ is strictly
positive on $\mathbb R$.
The two evaluations at $x=\log p/(2\pi)$
and $x=-\log p/(2\pi)$ sum to $(\sqrt\pi/\varepsilon)
(1+e^{-(\log p)^2/\varepsilon^2})$,
with the reflected term exponentially suppressed: at
$\varepsilon=0.05$ and $p=2$ it is below $10^{-82}$, and decreases
further for larger~$p$.
Collecting the dominant contribution,
\[
\mathrm{Main}_p(\varepsilon) \;=\; -\frac{1}{2\pi}\cdot\frac{\log p}{\sqrt p}
\cdot\frac{\sqrt\pi}{\varepsilon}
\bigl(1 + e^{-(\log p)^2/\varepsilon^2}\bigr)
= -\frac{\log p}{2\sqrt\pi\,\varepsilon\,\sqrt p}
\bigl(1 + e^{-(\log p)^2/\varepsilon^2}\bigr).
\]
As $\varepsilon\to 0$, the factor $e^{-(\log p)^2/\varepsilon^2}\to 0$,
giving $|\mathrm{Main}_p(\varepsilon)|\sim
(\log p)/(2\sqrt\pi\,\varepsilon\,\sqrt p)$.
The sign is negative because all factors in the leading expression are
positive and enter with a minus sign.
\end{proof}

\begin{corollary}[Main-term scaling; proved]
\label{cor:main-scaling}
As $\varepsilon\to 0$,
\[
|\mathrm{Main}_p(\varepsilon)|\sim
\frac{\log p}{2\sqrt\pi\,\varepsilon\sqrt p},
\]
giving a leading-order scaling of $1/\varepsilon$ with magnitude
constant $1/(2\sqrt\pi)\approx 0.2821$;
the signed main term carries an additional minus sign from the prime-sum formula.
\end{corollary}

At the reference parameter $\varepsilon=0.05$ and for $2\le p\le 53$,
the main term takes values in the interval $[-4.15,-2.77]$, minimised
at $p=7$ and maximised at $p=2$.

% -------------------------------------------------------
\section{Error Terms}
\label{sec:errors}

We now treat the three remaining contributions other than the constant
term.

\subsection{The gamma contribution}

\begin{theorem}[Gamma term is subleading; numerical]
\label{thm:gamma}
On the tested finite grid
$p\le 53$, $\varepsilon\in\{0.005,0.010,\ldots,0.100\}$,
the ratio $|\Gamma_p(\varepsilon)|/|\mathrm{Main}_p(\varepsilon)|$
is numerically consistent with linear scaling in $\varepsilon$;
a log--log regression gives slope~$1.001$.
At $\varepsilon=0.05$, the ratio lies in $[0.017,\,0.172]$
for all tested primes ($p\le 53$); in particular
$|\Gamma_p(0.05)|<|\mathrm{Main}_p(0.05)|$ for all $16$ primes $p\le 53$.
\end{theorem}

\begin{proof}[Numerical verification]
For every fixed $\varepsilon>0$, the Gaussian factor $e^{-\varepsilon^2 t^2}$
makes the integrand absolutely integrable on $\mathbb{R}$.
The claimed ratio range and log--log slope are established by direct
numerical quadrature over the grid
$p\le 53$, $\varepsilon\in\{0.005,0.010,\ldots,0.100\}$
(step $0.005$, giving $20$ values), pooling all $16$ primes;
the slope is the least-squares slope of
$\log(|\Gamma_p|/|\mathrm{Main}_p|)$ against $\log\varepsilon$.
No limit $\varepsilon\to 0$
is asserted. Establishing an Abel-regularised bound uniform in~$p$,
which would require explicit control of the oscillatory integral
as $\varepsilon\downarrow 0$, is part of Problem~\ref{op:bias-asymp}.
\end{proof}

\subsection{The prime-power and other-prime errors}

\begin{theorem}[Other-prime error is non-positive; proved]
\label{thm:other}
For every fixed prime $p\le\kappa$ and every $\varepsilon>0$,
\[
\mathrm{Err}_{\mathrm{other},p}(\varepsilon) \;\le\; 0.
\]
\end{theorem}

\begin{proof}
By the standard form of the explicit formula,
\[
\mathrm{Err}_{\mathrm{other}}(\varepsilon) \;=\; -\frac{1}{2\pi}
\sum_{\substack{q\le\kappa\\ q\neq p\\ q\text{ prime}}}
\frac{\log q}{\sqrt q}
\!\left[\widehat{h}_{p,\varepsilon}\!\left(\tfrac{\log q}{2\pi}\right)
      +\widehat{h}_{p,\varepsilon}\!\left(-\tfrac{\log q}{2\pi}\right)
\right]
+ (\text{higher-prime-power terms}).
\]
Each summand has $\log q>0$, $\sqrt q>0$, and
$\widehat{h}_{p,\varepsilon}(\cdot)>0$ on $\mathbb R$, with an overall
minus sign.
The higher-prime-power terms share the same sign structure with
$\Lambda(q^k)=\log q$ and the same positive transform.
Hence the full sum is non-positive.
The identical sign argument applies to all primes $q>\kappa$ and
higher prime powers with $q\neq p$: the positivity of
$\widehat h_{p,\varepsilon}$ on~$\mathbb R$ and the negativity of
the overall sign extend without restriction to the full prime-power range
in the explicit formula.
Hence $\mathrm{Err}_{\mathrm{other}}(\varepsilon)\le 0$ for the
complete GW sum.
\end{proof}

Theorem~\ref{thm:other} is a structural identity rather than an error
bound: the contribution from primes $q\neq p$ actively reinforces the
negativity of $\widetilde Z_p$.

\begin{theorem}[Truncation error is negligible; numerical]
\label{thm:trunc}
\[
\frac{|R_{p,N}(\varepsilon)|}{|\mathrm{Main}_p(\varepsilon)|}
\;\le\; 3\times 10^{-61}
\]
at the reference parameters $(\varepsilon,N)=(0.05,100)$ for every
prime $p\le 53$, where $R_{p,N}(\varepsilon)=\sum_{k>N}
e^{-\varepsilon^2\gamma_k^2}\cos(\gamma_k\log p)$ is the tail of the
zero-ordinate sum.
\end{theorem}

\begin{proof}
The triangle inequality gives
$|R_{p,N}(\varepsilon)|\le\sum_{k>N} e^{-\varepsilon^2\gamma_k^2}$.
We bound this tail by partial summation against the Riemann--von
Mangoldt density estimate $N(t+1)-N(t)\ll\log(t+2)$
(\cite[Thm.\,5.8]{IwaniecKowalski04},
with explicit zero-counting bounds from Bellotti--Wong~\cite{BellottiWong24}):
\[
\sum_{\gamma>\gamma_{100}} e^{-\varepsilon^2\gamma^2}
\;\ll\;
  e^{-\varepsilon^2 \gamma_{100}^2}\log(\gamma_{100}+2)
  + 2\varepsilon^2\int_{\gamma_{100}}^\infty e^{-\varepsilon^2 t^2}
    t\,\log(t+2)\,dt.
\]
At $(\varepsilon,N)=(0.05,100)$, $\gamma_{100}=236.524$, the Gaussian
factor $e^{-\varepsilon^2 \gamma_{100}^2}=e^{-0.0025\cdot 236.524^2}
\approx e^{-140}$ renders both terms negligible.
A standard Gaussian tail bound gives
$\int_{236}^\infty e^{-0.0025 t^2}\log t\,dt\ll e^{-140}$,
and the lower bound
$|\mathrm{Main}_p(\varepsilon)|\ge
\log 2/(2\sqrt{\pi}\cdot 0.05\cdot\sqrt{53})$
over the tested prime range yields
$|R_{p,N}|/|\mathrm{Main}_p|\le 3\times10^{-61}$.
Accurate evaluation of the partial sums follows
Brent--Platt--Trudgian~\cite{BrentPlattTrudgian21}.
For the $\gamma_k^2$-weighted tail: the function
$t^2 e^{-\varepsilon^2 t^2}$ is monotone decreasing for $t>1/\varepsilon=20$,
and $\gamma_{100}=236.524\gg 1/\varepsilon$.
Hence for $k>N$, $\gamma_k>\gamma_{100}$ and
$\gamma_k^2 e^{-\varepsilon^2\gamma_k^2}
 \le \gamma_{100}^2 e^{-\varepsilon^2\gamma_{100}^2}
 \approx \gamma_{100}^2\cdot e^{-140}$.
By partial summation the sum satisfies
\[
\sum_{k>N}\gamma_k^2 e^{-\varepsilon^2\gamma_k^2}
\;\ll\;
\gamma_{100}^2\,e^{-\varepsilon^2\gamma_{100}^2}\log(\gamma_{100}+2)
+\int_{\gamma_{100}}^\infty(2t+2\varepsilon^2 t^3)e^{-\varepsilon^2 t^2}
  \log(t+2)\,dt
\;\ll\;
\gamma_{100}^2\cdot e^{-140},
\]
which is negligible relative to $|B|$ at the reference parameters.
The numerical ratio $|R_{p,N}|/|\mathrm{Main}_p|\le 3\times10^{-61}$
is verified directly by \texttt{mpmath} computation at 50-digit precision
in the accompanying script~\cite{Code}: the sum
$\sum_{k=101}^{1000}e^{-\varepsilon^2\gamma_k^2}$
is evaluated numerically, confirming the stated bound for
all primes $p\le 53$ at the reference parameters.
The remaining infinite tail $k>1000$ is bounded analytically.
At $\gamma_{1000}\approx 1419$ and $\varepsilon=0.05$, the
Riemann--von Mangoldt estimate gives
\[
\sum_{k>1000}e^{-\varepsilon^2\gamma_k^2}
\;\ll\;
e^{-\varepsilon^2\gamma_{1000}^2}\log(\gamma_{1000}+2)
\;\approx\;
e^{-5034}\cdot 7.3
\;\approx\; 10^{-2184},
\]
which is negligible relative to the partial sum above and to the
stated bound $3\times10^{-61}$.
The full infinite-tail control used for the numerical verification
is obtained by combining the
numerically computed finite partial sum $(k=101,\ldots,1000)$
with this explicit analytic estimate.
\end{proof}

For the ordinate proxy $Z_{p,N}^+$ at the reference parameters,
the $N=100$ unweighted truncation tail is negligible relative to
$|\mathrm{Main}_p|$ (Theorem~\ref{thm:trunc}).
The analogous $\gamma_k^2$-weighted tail entering $B$ is
controlled by the same Gaussian mechanism; no separate quantified
bound is stated here, as the sign of $B=-19\,342.5$ is unaffected
by the tail at the stated numerical precision.

\begin{remark}[Asymptotic constant-term ratio; numerical]
\label{rem:rpinf}
The proxy ratio $r^+_{p,N}(\varepsilon):=
|\mathrm{Const}^+_{p,N}(\varepsilon)|/|\mathrm{Main}_p(\varepsilon)|$
is observed, on the smallest tested $\varepsilon$-values, to be
consistent with decrease toward a limiting value as $\varepsilon\downarrow 0$.
At the smallest tested values ($\varepsilon\le 0.015$), all ratios
lie below~$0.82$ for every prime $p\le 53$, consistent with convergence
to a limit $r_p^{\infty}\le\tfrac12$ as $\varepsilon\to 0$.
An analytic proof of $r_p^{\infty}=\tfrac12$ requires an explicit
asymptotic expansion of $\mathrm{Const}_p(\varepsilon)$ as $\varepsilon\to 0$
beyond the leading $e^{\varepsilon^2/4}\cosh(\tfrac{\log p}{2})$ pole term
evaluated in \S\,\ref{sec:gw}; this is part of Problem~\ref{op:bias-asymp}.
On the smallest tested values, the data are consistent with
$\mathrm{Const}_p(\varepsilon)+\mathrm{Main}_p(\varepsilon)$ decreasing
toward zero; no asymptotic proof is supplied here.
\end{remark}

% -------------------------------------------------------
\section{The Constant Term and the Sign-Crossover Localisation}
\label{sec:const}

The constant term $\mathrm{Const}_p(\varepsilon)$ is the only contribution
whose full asymptotic structure is not explicitly resolved in the present work.
It collects the
$s=0,1$ endpoint contributions in the completed Guinand--Weil formula
(associated with the zeta singularity at $s=1$ and its reflection
at $s=0$ under the functional equation), the $\log(2\pi)$ normalisation,
and the remaining endpoint and normalisation terms not included in
the displayed prime-power sums.
It is the only term of the decomposition that resists a short analytic
treatment.
We state a numerical result on its asymptotic ratio to the main term
and, independently, a numerical result localising the sign-crossover
of $Z_{p,N}^+(\varepsilon)$ on a finite $\varepsilon$-grid.

\begin{theorem}[Finite-grid proxy-ratio pattern and sign-crossover localisation; numerical]
\label{thm:ratio-signcross}
Numerically, the following hold.

\emph{(a)} The proxy ratio
$r^+_{p,N}(\varepsilon):=|\mathrm{Const}^+_{p,N}(\varepsilon)|/
|\mathrm{Main}_p(\varepsilon)|$ appears unimodal on the tested
$\varepsilon$-grid, with an observed maximum near the largest tested
value $\varepsilon=0.050$;
no continuous unimodality statement is proved.
The maximum exceeds~$1$ for $p=37$ and $p=53$
(attaining $1.17$ and $1.21$ respectively at $\varepsilon=0.05$).
On the smallest tested values $\varepsilon\in\{0.005,0.010,0.015\}$,
all ratios lie below~$0.82$ for every prime $p\le 53$.
The numerical data are consistent with decrease toward
$r_p^{\infty}\le\tfrac12$ for the underlying analytical object;
no analytic limit is proved
(Remark~\ref{rem:rpinf}).

\emph{(b)} On the sampled grid, the bias changes sign between
the tested values $\varepsilon=0.020$ and $\varepsilon=0.025$.
At the tested value $\varepsilon=0.020$, one finds
\[
Z_{p,N}^+(\varepsilon) \;<\; 0 \qquad
\text{for every prime } p\le 53
\quad\text{(ordinate proxy } Z_{p,N}^+ = \Re Z_{p,\varepsilon,N}\text{)}.
\]

Both statements are numerical and refer to the cutoff $\kappa=53$.
The location of the sign-crossover and the sign at smaller tested
values of $\varepsilon$ are empirical findings on a finite grid; no
continuous interval-persistence statement is asserted.
The analytical question behind (a) is documented as
Problem~\ref{op:bias-asymp} of~\S\,\ref{sec:open}.
\end{theorem}

\emph{Discussion.}
The $s=1$ endpoint evaluation gives
$h^{\mathrm{even}}_{p,\varepsilon}(i/2)+h^{\mathrm{even}}_{p,\varepsilon}(-i/2)
= e^{\varepsilon^2/4}(\sqrt p+1/\sqrt p)$,
which is $O(1)$ as $\varepsilon\downarrow0$ for fixed $p$
(see Edwards~\cite[\S1.8, \S1.16, and \S3.6]{Edwards}).
The observed $\varepsilon^{-1}$ component of the proxy residual
$\mathrm{Const}^+_{p,N}$ arises from the full
residual in the chosen partition and is not an analytic property
of the endpoint term alone; an explicit derivation is part of
Problem~\ref{op:bias-asymp}.
The limiting coefficient of the proxy ratio is not produced by the
endpoint term alone.
It arises from the positive-ordinate proxy normalisation: under the
GW bridge, $Z_{p,N}^+\approx\tfrac12\widetilde Z_p^{\mathrm{full}}$,
and hence
$\mathrm{Const}^+_{p,N}\approx
\tfrac12\mathrm{Const}_p-\tfrac12(\mathrm{Main}_p+\Gamma_p
+\mathrm{Err}_{p^k}+\mathrm{Err}_{\mathrm{other}})$.
Since $\mathrm{Main}_p(\varepsilon)\sim
-\frac{\log p}{2\sqrt\pi\sqrt p}\cdot\varepsilon^{-1}$, this
identity explains the observed leading contribution
$\mathrm{Const}^+_{p,N}\approx-\tfrac12\mathrm{Main}_p$ and hence
the expected ratio $r_p^\infty\to\tfrac12$, modulo the open GW bridge
and analytic control of the remaining terms.
Numerical evaluation of $\mathrm{Const}^+_{p,N}(\varepsilon)$ for
$\varepsilon\in\{0.005,0.010,0.015,0.020,0.025,0.030,0.050\}$ across
all primes $p\le 53$ shows that the proxy ratio
$|\mathrm{Const}^+_{p,N}(\varepsilon)|/|\mathrm{Main}_p(\varepsilon)|$ is a
function whose values on the tested grid exhibit a single peak near
$\varepsilon=0.050$ and limiting value
below~$0.82$ on the smallest tested grid values ($\varepsilon\le 0.015$),
consistent with decrease toward a limiting value
$r_p^{\infty}\le\tfrac12$ (Remark~\ref{rem:rpinf}).
For $p\in\{37,53\}$ this maximum narrowly exceeds unity (attaining
$1.183$ and $1.216$ respectively), which produces two pointwise
exceptions to the bias at $\varepsilon=0.05$; see below.
The sign-crossover localisation in (b) is obtained from the grid
computation directly, not as an analytic consequence of (a).

\emph{Numerical evidence.}
At $\varepsilon=0.05$ and $\kappa=53$, the proxy-residual values satisfy
$\mathrm{Const}^+_{p,N}\in[2.44,3.93]$ with mean $3.27$ and standard
deviation $0.46$ (coefficient of variation $14\%$).
The displayed $\mathrm{Const}^+_{p,N}(\varepsilon)$ values are computed
as the residual
$\mathrm{Const}^+_{p,N}(\varepsilon) := Z_{p,N}^+(\varepsilon)
  - \mathrm{Main}_p(\varepsilon)
  - \Gamma_p(\varepsilon)
  - \mathrm{Err}_{p^k}(\varepsilon)
  - \mathrm{Err}_{\mathrm{other}}(\varepsilon)$
at the reference parameters $(\varepsilon,N)=(0.05,100)$.
These numerical values are the ordinate-proxy residuals: the
left-hand side is the positive-ordinate half-sum $Z_{p,N}^+$,
while the subtracted terms ($\mathrm{Main}_p,\Gamma_p,
\mathrm{Err}_{p^k},\mathrm{Err}_{\mathrm{other}}$)
are the full Guinand--Weil terms of Proposition~\ref{prop:gw}.
Since $Z_{p,N}^+\approx\tfrac12\widetilde Z_p$ under the bridge,
a short computation gives
\[
\mathrm{Const}^+_{p,N}(\varepsilon)
\;\approx\;
\tfrac12\mathrm{Const}_p(\varepsilon)
-\tfrac12\bigl(\mathrm{Main}_p(\varepsilon)
+\Gamma_p(\varepsilon)+\mathrm{Err}(\varepsilon)\bigr),
\]
so $\mathrm{Const}^+_{p,N}$ is not a half-normalisation of
$\mathrm{Const}_p$; the two differ by the half-sum of the
other terms.
The proxy ratio
$r^+_{p,N}(\varepsilon)=|\mathrm{Const}^+_{p,N}|/|\mathrm{Main}_p|$
approaches a value $\le\tfrac12$ as $\varepsilon\to 0$
(Remark~\ref{rem:rpinf}) because $|\mathrm{Main}_p|\sim\varepsilon^{-1}$
dominates the $O(1)$ difference term in the limit.
No exact identity with $\mathrm{Const}_p$ of
Proposition~\ref{prop:gw} is claimed without the bridge.
Representative values are listed below.

\begin{center}
\begin{tabular}{cccc}
\toprule
$p$ & $\mathrm{Main}_p(0.05)$ & $\Gamma_p(0.05)$
    & $\mathrm{Const}^+_{p,N}(0.05)$ \\
\midrule
$2$  & $-2.77$ & $-0.475$ & $2.44$ \\
$7$  & $-4.15$ & $-0.193$ & $3.47$ \\
$23$ & $-3.69$ & $-0.105$ & $3.79$ \\
$37$ & $-3.35$ & $-0.082$ & $3.93$ \\
$53$ & $-3.08$ & $-0.069$ & $3.74$ \\
\bottomrule
\end{tabular}
\end{center}

At $\varepsilon=0.05$, the pointwise inequality
$Z_{p,N}^+(0.05)<0$ holds for $14$ of the $16$ primes
$p\le 53$.
The exceptions are $p=37$ and $p=53$, for which
$Z_{p,N}^+(0.05)\approx+0.50$ and $+0.59$ respectively.
This is \emph{not} a failure of the underlying decomposition: it
reflects the finite-$\varepsilon$ maximum of the ratio $r(\varepsilon)$
near $\varepsilon\approx 0.05$, which narrowly exceeds unity for
exactly these two primes.
At the tested value $\varepsilon=0.020$ the pointwise bias holds for
all $16$ primes without exception.
At $\varepsilon=0.025$, only $p=53$ fails the pointwise negativity
condition ($Z_{53,N}^+(0.025)=+0.389>0$), while $p=37$ remains
negative ($Z_{37,N}^+(0.025)=-0.356<0$).
At $\varepsilon=0.030$, both exceptions $p=37,53$ are present.
The sign-crossover for $p=53$ lies in $(0.020,\,0.025)$, and
that for $p=37$ lies in $(0.025,\,0.030)$.
Whether the bias persists continuously throughout
$\varepsilon\in(0,0.020]$ is not established here.

No structural explanation for the exceptions $p=37,53$ at
$\varepsilon=0.05$ is currently available.
The phenomenon appears to be a finite-$\varepsilon$, finite-$\kappa$
manifestation of the observed unimodal peak of $r(\varepsilon)$ near
$\varepsilon=0.050$, which for exactly these two
primes narrowly exceeds unity.
Whether this reflects an underlying arithmetic structure or is a
numerical coincidence at $\kappa=53$ is not addressed here.

The observed negativity of $Z_{p,N}^+(\varepsilon)$ should
be read neither as a prime race in the sense of
Rubinstein--Sarnak~\cite{RubinsteinSarnak} nor as a Landau--Gonek
prime-power asymptotic; it is a different, prime-indexed sign
statistic for Gaussian-smoothed explicit-formula evaluations.

\emph{Status.}
The tested values of $r^+_{p,N}(\varepsilon)$ are numerically consistent
with decrease toward a value satisfying $r_p^{\infty}\le\tfrac12$
(Remark~\ref{rem:rpinf}).
The sign of the ordinate proxy $Z_{p,N}^+$ remains the
subject of Problem~\ref{op:bias-asymp}.
Establishing a uniform bound on the finite-$\varepsilon$ correction,
which would close the asymptotic bias conjecture in its sharpest form,
is part of Problem~\ref{op:bias-asymp} of~\S\,\ref{sec:open}.
The empirical sign-crossover localisation above, at $\kappa=53$, is
the sharpest statement that the present methods allow.


\begin{remark}[Relation to Balanzario--C\'ardenas Romero]
Balanzario and C\'ardenas Romero~\cite{BalanzarioCardenas25}
study a Gaussian-Hermite-weighted explicit formula
$S(\xi,k) = (1/\xi)\sqrt{k/2\pi}\,
\sum_n \exp\{-2(k/\xi^2)(\gamma_n - \xi)^2\}$
with a negative prime-side leading term of the form
$-c_2\,\Lambda(n)/\sqrt{n}\cdot\exp\{-(\xi^2/4k)\log^2 n\}$,
which is structurally analogous to our
$\mathrm{Main}_p(\varepsilon) = -\log p/(2\sqrt{\pi}\,\varepsilon\sqrt{p})$.
Their parametrisation uses Hermite polynomials and prime-power
sums $\Lambda(n)$; our single-prime, leading-$\varepsilon$ form
with explicit constant $1/(2\sqrt{\pi})$ appears to be a new
observation.
\end{remark}
% -------------------------------------------------------
\section{The Integrated Bias}
\label{sec:intbias}

The two expressions for $B$ used below are the same quantity in two
representations: the primal sum
\[
B \;=\; \sum_{k,p} e^{-\varepsilon^2\gamma_k^2}\,(\gamma_k\log p)^2\,
         \cos(\gamma_k\log p)
\]
and the dual sum $B=\sum_p(\log p)^2\,\mathrm{Re}\,Z_p^{(2)}$.
The equivalence between the two is the Curvature-Bias Decomposition
(cf.~\cite{Paper5}, \S5)
and is treated here as imported.
All numerical values in this section refer to the same object.

We use the dual representation as a starting point.
The curvature bias at $\sigma=\tfrac12$ is encoded in the weighted sum
\[
B \;:=\; \sum_{p\le\kappa}(\log p)^2\,\mathrm{Re}\,Z_p^{(2)},
\]
which was introduced in~\cite{Paper5} as $-\tfrac12$ times the
$\sigma=\tfrac12$ evaluation of the second $\sigma$-derivative of the
trace identity, via the identity $O''(\tfrac12)=-2B$.
Its structural decomposition is algebraic and does not depend on any
unproved hypothesis.

\begin{proposition}[Curvature--bias identity; proved]
\label{prop:curv-bias}
Let $O(\sigma)=\tfrac12\sum_{k,p} e^{-\varepsilon^2\gamma_k^2}
\cos(2\sigma\gamma_k\log p)$ and define
\[
B \;:=\; \sum_{k,p} e^{-\varepsilon^2\gamma_k^2}\,(\gamma_k\log p)^2\,
         \cos(\gamma_k\log p).
\]
Then
\[
O''(\tfrac12) \;=\; -2\,B.
\]
\end{proposition}

\begin{proof}
Direct differentiation gives
\[
O'(\sigma) \;=\; -\sum_{k,p} e^{-\varepsilon^2\gamma_k^2}\,\gamma_k\log p\,
                \sin(2\sigma\gamma_k\log p),
\]
\[
O''(\sigma) \;=\; -2\sum_{k,p} e^{-\varepsilon^2\gamma_k^2}\,
                (\gamma_k\log p)^2\,\cos(2\sigma\gamma_k\log p).
\]
Evaluation at $\sigma=\tfrac12$ yields
\[
O''(\tfrac12) \;=\; -2\sum_{k,p} e^{-\varepsilon^2\gamma_k^2}\,
                (\gamma_k\log p)^2\,\cos(\gamma_k\log p) \;=\; -2\,B.
\qedhere
\]
\end{proof}

\begin{remark}
This identity shows that all curvature information at $\sigma=\tfrac12$
is encoded in a single weighted cosine interaction between primes and zeros.
\end{remark}

Numerically, at reference parameters $(\kappa,\varepsilon,N)=(53,0.05,100)$,
the direct curvature sum equals $B=-19\,342.5$, hence
$O''(\tfrac12)=-2B=+38\,685$.
This value coincides with $\mathrm{Tr}(L)$, where
$L=-d^2\Tt(\sigma)/d\sigma^2\big|_{\sigma=1/2}$ is the curvature
operator of~\cite{Paper5}, and provides an independent closed-form
cross-check for the identity above.

\begin{remark}
The quantity $B$ is defined as in~\cite[\S5]{Paper5}, and the identity
$O''(\tfrac12)=-2B$ follows from the explicit form of $O(\sigma)$
given there.
The numerical value $B=-19\,342.5$ at $(\kappa,\varepsilon,N)=(53,0.05,100)$
is the one quoted throughout the series.
\end{remark}

\begin{theorem}[Integrated proxy bias is negative; numerical]
\label{thm:intbias}
At $(\kappa,\varepsilon,N)=(53,0.05,100)$,
\[
B_{\mathrm{int}}^+(0.05,100) \;:=\; \sum_{p\le 53}(\log p)^2\,
Z_{p,N}^+(0.05) \;=\; -42.21 \;<\; 0.
\]
\end{theorem}

\begin{observation}[Direct curvature sum; numerical]
\label{obs:B}
At $(\kappa,\varepsilon,N)=(53,0.05,100)$,
$B=-19\,342.5<0$, equivalently $O''(\tfrac12)=-2B=+38\,685>0$.
\end{observation}

When the parameters $(\varepsilon,N)=(0.05,100)$ are fixed throughout,
we write $B_{\mathrm{int}}$ as shorthand for $B^+_{\mathrm{int}}(0.05,100)$.

\begin{remark}
The values $B$ and $B_{\mathrm{int}}$ refer to two \emph{different}
quantities, not to a single quantity in two forms.
$B$ is the direct curvature sum weighted by $(\gamma_k\log p)^2$,
obtained by differentiating the trace identity and evaluating at
$\sigma=\tfrac12$; $B_{\mathrm{int}}^+(\varepsilon,N)$ is the integrated proxy bias,
weighted only by $(\log p)^2$ and summed over $Z_{p,N}^+(\varepsilon)$.
The two differ by a reweighting through the $\gamma_k^2$-factors
absent in the second.
Both are numerically negative at the reference parameters, and the
shared negative sign is a numerical finding rather than the consequence
of a proved equivalence.
Proposition~\ref{prop:curv-bias} connects $B$ --- not $B_{\mathrm{int}}$
--- to the curvature $O''(\tfrac12)$.
Whether pointwise negativity of $Z_{p,N}^+(\varepsilon)$, or
the integrated proxy bias $B_{\mathrm{int}}^+(\varepsilon,N)<0$, implies the
corresponding negativity of $B$ at the same parameters is an open
question; it is formulated as Problem~\ref{op:transfer} of~\S\,\ref{sec:open}.
\end{remark}

\begin{remark}
Although $Z_{p,N}^+(0.05)>0$ at $p=37$ and $p=53$, the
$(\log p)^2$-weighted sum $B_{\mathrm{int}}$ remains decisively
negative: the negative contributions from the remaining $14$ primes
dominate the positive contributions from these two.
Concretely, $(\log 37)^2\cdot 0.50+(\log 53)^2\cdot 0.59\approx
6.52+9.30=15.82$, against a cumulative negative contribution of
$-58.03$ from the other primes, yielding the net value $-42.21$.
The integrated bias is robust against individual-prime deviations as
long as these remain isolated.
\end{remark}

% -------------------------------------------------------
\section{Positive Curvature at the Critical Line}
\label{sec:curvcor}

We now record the curvature statement that follows from
Observation~\ref{obs:B} and Proposition~\ref{prop:curv-bias}, and,
separately, the conditional strict-local-extremum statement that
follows once stationarity is added as a hypothesis.

\begin{corollary}[Positive curvature at $\sigma=\tfrac12$; numerical]
\label{cor:curvature}
At the reference parameters $(\kappa,\varepsilon,N)=(53,0.05,100)$,
the direct curvature sum $B$ is numerically negative
(Observation~\ref{obs:B}), $B=-19\,342.5<0$.
By Proposition~\ref{prop:curv-bias}, this is equivalent to
\[
O''(\tfrac12) \;=\; -2B \;=\; +38\,685 \;>\; 0.
\]
\end{corollary}

\begin{proof}
By Observation~\ref{obs:B}, $B=-19\,342.5$.
By Proposition~\ref{prop:curv-bias}, $O''(\tfrac12)=-2B$.
\end{proof}

\begin{remark}
If, in addition to $B<0$ at the reference parameters (here numerically
established), the stationarity condition $O'(\tfrac12)=0$ holds at the
same parameters, then $\sigma=\tfrac12$ is a strict local minimum of
$O$, and the trace
\[
\mathrm{Tr}\,\Tt(\sigma) \;=\; \DSEL - O(\sigma)
\]
attains a strict local maximum at $\sigma=\tfrac12$.

The exact vanishing condition $O'(\tfrac12)=0$ is stronger than the
first-derivative cancellation bound formulated in Problem~\ref{op:stat}
and remains open.
Numerically, $O'(\tfrac12)=+2.475$ is approximately $4\%$ of the
natural scale $W_{\varepsilon,N}\cdot P(\kappa)$, consistent with the cancellation expected
from the odd kernel $\sin(\gamma_k\log p)$, but no analytical bound on
$|O'(\tfrac12)|$ is established in the present paper.
\end{remark}

\begin{remark}
The paper~\cite{Paper5} recorded four further numerical signatures
concentrating at $\sigma=\tfrac12$: a $+47\%$ trace spike relative to
the $\sigma$-neighbourhood average, the dominance of the leading
eigenvalue $\mu_1$, the minimum of the spectral gap, and the joint
maximum of $\mathrm{Tr}(L)$ and $\lambda_{\max}(L)$.
Each of these is numerically robust and independent of the bias
conditions established here.
They constitute additional numerical signatures from~\cite{Paper5} at
the same reference parameters, but do not enter the proof of
Corollary~\ref{cor:curvature} or of the conditional remark above.
\end{remark}

\begin{remark}
The curvature statement identifies $\sigma=\tfrac12$ as a geometrically
distinguished point of the spectral trace at the reference parameters.
It is a structural ingredient in a possible route to Weil positivity
in the Lagarias formulation~\cite{Lagarias}
(Problem~\ref{op:weil-bridge} of~\S\,\ref{sec:open}), but does not by
itself establish it.
No implication concerning the Riemann Hypothesis is asserted or
implied by Corollary~\ref{cor:curvature}, nor by the conditional remark
above when the stationarity hypothesis is imposed.
The curvature identity therefore provides a canonical second-order
signature of the critical line, independent of global hypotheses.
\end{remark}

% -------------------------------------------------------
\section{Open Problems}
\label{sec:open}

The present work reduces the question of curvature at the critical
line to five well-posed subproblems.
We name each and describe the tools we expect to be relevant.
The problems are listed in order of what we consider tractability.

\begin{openproblem}[First-derivative cancellation and localisation]\label{op:stat}
The trivial triangle inequality gives
\[
|O'_{(\kappa,\varepsilon,N)}(\tfrac12)|
\;\le\;
P(\kappa)\sum_{k=1}^N\gamma_k e^{-\varepsilon^2\gamma_k^2}
\;=\;
\widetilde C_{\varepsilon,N}\,W_{\varepsilon,N}P(\kappa),
\qquad
\widetilde C_{\varepsilon,N}
:=\frac{\sum_{k=1}^N\gamma_k e^{-\varepsilon^2\gamma_k^2}}
       {\sum_{k=1}^N e^{-\varepsilon^2\gamma_k^2}},
\]
which is finite for fixed $(\varepsilon,N)$ but does not represent
cancellation. The nontrivial problem is to prove a first-derivative
\emph{cancellation estimate} of the form: for some $C_0>0$ and $A\ge1$,
\[
\bigl|S_\kappa(\gamma)\bigr|
\;\le\; C_0\,\frac{P(\kappa)}{\gamma(\log\gamma)^A},
\qquad
S_\kappa(\gamma):=\sum_{p\le\kappa}(\log p)\sin(\gamma\log p),
\]
uniformly in $\kappa\ge 2$ and $\gamma\ge\gamma_1$. This would imply
\[
|O'_{(\kappa,\varepsilon,N)}(\tfrac12)|
\;\le\;
\frac{C_0\,W_{\varepsilon,N}\,P(\kappa)}{(\log\gamma_1)^A},
\]
which is a genuine logarithmic saving over the trivial bound.
The numerical value at reference parameters is $O'(\tfrac12)=+2.475$,
approximately $4\%$ of $W_{\varepsilon,N}\cdot P(\kappa)$, consistent with the cancellation
expected from the odd kernel $\sin(\gamma_k\log p)$.
The expected tools are: equidistribution of the sequence
$\{\gamma_k\log p\bmod 2\pi\}$ under appropriate averaging; a discrete
Weyl-type bound exploiting the independent behaviour of zero heights
and prime logarithms.
At the reference parameters, $B<0$ is numerically established
(Corollary~\ref{cor:curvature}); the remaining open hypothesis is the
prime-sum cancellation estimate $|S_\kappa(\gamma)|\le C_0 P(\kappa)/(\gamma(\log\gamma)^A)$
stated above.
Such a bound would become a localisation statement only when combined
with a local convexity lower bound and a smallness condition relative
to that lower bound, as detailed below.
Specifically, if there exist $\delta>0$ and $m>0$ such that
$O''(\sigma)\ge m$ on $[\tfrac12-\delta,\tfrac12+\delta]$
and $|O'(\tfrac12)|<m\delta$, then there exists
$\sigma_*\in[\tfrac12-\delta,\tfrac12+\delta]$ with
$O'(\sigma_*)=0$ and
\[
|\sigma_* - \tfrac12| \;\le\; \frac{|O'(\tfrac12)|}{m}.
\]
The reference-parameter estimate gives
$\sigma_* - \tfrac12
\approx -O'(\tfrac12)/O''(\tfrac12) \approx -6.4\times10^{-5}$.
An exact local minimum at $\sigma=\tfrac12$ requires additionally
$O'(\tfrac12)=0$; this stronger condition remains open.
\end{openproblem}

\begin{remark}[Possible conditional route to stationarity]
\label{rem:stat-conditional}
A possible partial route toward Problem~\ref{op:stat} would be a cancellation
estimate for the prime exponential sum
\[
S_\kappa(\gamma):=\sum_{p\le\kappa}(\log p)\sin(\gamma\log p).
\]
Since
\[
O'(\tfrac12)=-\sum_k e^{-\varepsilon^2\gamma_k^2}\gamma_k\,S_\kappa(\gamma_k),
\]
an estimate of the form
\[
|S_\kappa(\gamma)|\le C_0\,\frac{P(\kappa)}{\gamma(\log\gamma)^A},
\qquad P(\kappa):=\sum_{p\le\kappa}\log p,
\]
for constants $C_0>0$ and $A\ge 1$ (uniform in $\kappa\ge 2$ and
$\gamma\ge\gamma_1$) would imply
\[
|O'(\tfrac12)|\;\le\;
C_0 P(\kappa)\sum_k \frac{e^{-\varepsilon^2\gamma_k^2}}{(\log\gamma_k)^A}
\;\le\;
\frac{C_0\,W_{\varepsilon,N}\,P(\kappa)}{(\log\gamma_1)^A}.
\]
No such estimate is proved here. The hypothesis is analogous to
Vinogradov--Korobov-type cancellation (\cite{Bellotti24a,Bellotti24b},
\cite{MossinghoffTrudgianYang24}, \cite{Ramare12}) and is strictly
weaker than RH. The structural obstruction to an unconditional proof
is described in Remark~\ref{rem:selection-direct}.
Numerically, $O'(\tfrac12)=+2.475$ and
$W_{\varepsilon,N}\cdot P(\kappa)/\log\gamma_1\approx 23.3$
at the reference parameters $(\varepsilon,N)=(0.05,100)$,
consistent with the route above.
The route falls short of an unconditional proof: the logarithmic saving
required to establish $|S_\kappa(\gamma)|\le C_0 P(\kappa)/(\gamma(\log\gamma)^A)$
uniformly in $\kappa$ is not available; see Remark~\ref{rem:selection-direct}.
\end{remark}

\begin{remark}[Why non-vanishing on $\Re s=1$ is insufficient]
\label{rem:selection-direct}
Direct unconditional proofs of the stationarity bound face a
structural obstruction: the function $-\zeta'/\zeta(s-i\gamma)$
has a pole at $s=1+i\gamma$, so the non-vanishing $\zeta(1+it)\neq 0$
(Hadamard~\cite{Hadamard}, de~la~Vall\'ee~Poussin) alone does not prevent
cancellation failure
in the sum $\sum_{p\le\kappa}(\log p)\sin(\gamma\log p)$.
Thus non-vanishing on $\Re s=1$ alone is insufficient for this route.
A cancellation estimate of the strength displayed in
Remark~\ref{rem:stat-conditional}, or an alternative estimate of
comparable strength, is still required for this route.
\end{remark}

\begin{openproblem}[Asymptotic ordinate bias and finite-$N$ transfer]\label{op:bias}
\label{op:bias-asymp}

\emph{(a) Asymptotic target.}
Let $Z_p^{\infty}(\varepsilon):=
2\sum_{k=1}^{\infty}e^{-\varepsilon^2\gamma_k^2}\cos(\gamma_k\log p)$
be the infinite ordinate proxy. Prove analytically that
$Z_p^{\infty}(\varepsilon)<0$ for every prime $p\le\kappa$ and all
sufficiently small $\varepsilon>0$.

\emph{(b) Finite-$N$ transfer.}
For specified $(N,\varepsilon)$, establish
$|R_{p,N}^{\infty}(\varepsilon)|<|Z_p^{\infty}(\varepsilon)|$
uniformly for $p\le\kappa$, where
$R_{p,N}^{\infty}(\varepsilon):=2Z_{p,N}^+(\varepsilon)-Z_p^{\infty}(\varepsilon)$
is the tail beyond the $N$-th zero. Then $Z_{p,N}^+(\varepsilon)<0$ follows.

The numerical behaviour of $r^+_{p,N}(\varepsilon)$
(Remark~\ref{rem:rpinf})
suggests the expected sign for the leading asymptotic of the full GW object,
but the limit $r_p^{\infty}\le\tfrac12$ is not proved analytically;
the corresponding sign of $Z_p^{\infty}$ is not directly
implied (the GW bridge remains open; see Paper~5, Open Problem~6.5 (GW bridge)).
The open question is an analytic proof that $Z_p^{\infty}(\varepsilon)<0$
for all primes and small but positive~$\varepsilon$.
Numerically, pointwise negativity of $Z_{p,N}^+$ is established on the
tested grid at $\varepsilon=0.020$ for all primes $p\le 53$.
Closing this gap would prove the asymptotic pointwise-bias target in
part~(a) and, together with the finite-$N$ transfer in part~(b),
would give a theorem-level sign statement for $Z_{p,N}^+$ in the
corresponding small-$\varepsilon$ regime.
It would \emph{not} by itself prove the finite-grid unimodality
pattern or the observed grid transition of
Theorem~\ref{thm:ratio-signcross}; those remain numerical unless
separately analysed.
The expected tools are: explicit expansion of the $s=0,1$ endpoint
terms in the completed Guinand--Weil formula with the correct
$\cosh$-evaluation of the test function; a Riemann--Lebesgue-type bound on the digamma contribution
with explicit constant; a cancellation argument between the
$\varepsilon^{-1}$-coefficients of $\mathrm{Main}_p$ and
$\mathrm{Const}_p$.
The problem is formulated entirely within classical analytic number theory.
\end{openproblem}

\begin{openproblem}[Uniformity of positive curvature]\label{op:uniform}
\label{op:uniformity}
Determine whether the inequality $B(\varepsilon,\kappa)<0$, and hence
$O''(\tfrac12)>0$, persists at parameter points
$(\kappa,\varepsilon,N)$ other than the reference values
$(53,0.05,100)$ tested here.
Corollary~\ref{cor:curvature} establishes the inequality at a single
point of the parameter space.
Whether it extends to a neighbourhood, to a larger range of cutoffs
$\kappa$, or to a different choice of smoothing $\varepsilon$ is a
distinct question: the direct curvature sum depends non-trivially on
the $\gamma_k^2$-weights and on the oscillatory phases $\gamma_k\log p$,
and finite parameter changes could in principle reverse the sign.
The expected tools are: a systematic numerical survey across a
parameter grid; an asymptotic expansion of $B(\varepsilon,\kappa)$ for
large $\kappa$; and, if possible, an analytic reduction via the
primal-to-dual identity $B=\sum_p(\log p)^2\,\mathrm{Re}\,Z_p^{(2)}$.
\end{openproblem}

\begin{openproblem}[Integrated-to-direct transfer]\label{op:transfer}
Establish or disprove the implication
\[
\bigl(\,Z_{p,N}^+(\varepsilon)<0 \text{ for all primes }
p\le\kappa\,\bigr)
\;\Longrightarrow\;
B(\varepsilon,\kappa)<0
\]
at fixed $(\varepsilon,\kappa)$, where
\[
B(\varepsilon,\kappa) \;:=\; \sum_{k,p} e^{-\varepsilon^2\gamma_k^2}\,
(\gamma_k\log p)^2\,\cos(\gamma_k\log p).
\]
At the reference parameters $(\kappa,\varepsilon,N)=(53,0.05,100)$,
$B^+_{\mathrm{int}}$ and $B$ are numerically negative
(Theorem~\ref{thm:intbias}, Corollary~\ref{cor:curvature}),
while the pointwise inequality $Z_{p,N}^+(0.05)<0$ holds for
$14$ of the $16$ primes $p\le 53$, with exceptions $p=37,53$
(Theorem~\ref{thm:ratio-signcross}).
On the tested grid at $\varepsilon=0.020$, the pointwise inequality
holds for all primes $p\le 53$.
No analytical implication among these sign statements is established.
A proof of the displayed implication would extend
Corollary~\ref{cor:curvature} from a reference-parameter statement to
a parametrised statement valid wherever the pointwise hypothesis can
be verified.
The expected tools are: the explicit prime-side representation of $B$,
the uniform positivity of the Fourier transform
$\widehat{h}_{p,\varepsilon}$, and a rearrangement argument
distributing the $\gamma_k^2$-weights across the pointwise
contributions.
The analogous implication with $B^+_{\mathrm{int}}$ in place of $B$ is
immediate from the positivity of the weights $(\log p)^2$:
if $Z_{p,N}^+(\varepsilon)<0$ for all $p\le\kappa$, then
$B^+_{\mathrm{int}}=\sum_p(\log p)^2 Z_{p,N}^+<0$.
The nontrivial question is whether any sign information survives
the additional $\gamma_k^2$-reweighting in~$B$.
The problem is formulated entirely within classical analytic number theory.
\end{openproblem}

\begin{openproblem}[Lagarias--Weil bridge]\label{op:weil}
\label{op:weil-bridge}
The present result introduces a finite-cutoff curvature observable
for Gaussian-smoothed explicit-formula data; it should be compared
with Bombieri's finite truncations of Weil's quadratic
functional~\cite{Bombieri} and with the Li/Bombieri--Lagarias
positivity framework~\cite{Li,BombieriLagarias}, but it is not yet a
positivity theorem for the full Weil distribution.

Problem~\ref{op:weil-bridge} asks whether there exists a transfer
construction under which the finite-cutoff curvature condition
$O''(\tfrac12)=-2B>0$, together with stationarity
$O'(\tfrac12)=0$ and an appropriate limiting procedure, yields
Weil positivity for a specified Gaussian-generated subclass of
Lagarias-admissible test functions
(entire of order~2, explicit-formula admissibility conditions).
This problem is stated here as a named research target, not as a proved implication.
Establishing this route precisely would require specifying the
admissible test-function class, constructing an explicit transfer from
the finite-parameter curvature identity $O''(\tfrac12)=-2B$ to the
limiting Weil functional, and proving positivity on that limiting class;
none of these steps is supplied in the present paper.
Its resolution is beyond the methods of the present paper and requires
either a direct positivity argument on the Weil distribution or an
external route via Connes-type operator positivity.
We are not aware of a direct transfer theorem connecting the
curvature statement $O''(\tfrac12)>0$ or the Lagarias formulation
of Weil positivity~\cite{Lagarias} to Bombieri's quadratic
functional~\cite{Bombieri} or the Connes--Consani
Weil-positivity framework~\cite{ConnesConsani23}.
The two bodies of work remain disjoint as of the date of
this paper; Problem~\ref{op:weil-bridge} is therefore
genuinely open.
No implication toward the Riemann Hypothesis is asserted by the
present paper; the problem above is a named research target,
not an established claim.

Problems~\ref{op:stat}--\ref{op:transfer} are formulated entirely within
classical analytic number theory; the present problem is structurally
the principal outstanding question of the programme.
\end{openproblem}

\begin{remark}[Outlook: the Weil transfer operator]
\label{rem:weil-transfer-outlook}
The curvature identity $O''(\tfrac12)=-2B$ established in this
paper admits a natural continuation: a possible construction of a
finite-dimensional Weil transfer matrix $W_{\kappa,\varepsilon}$
whose entries would be evaluations of the Weil quadratic form on the
admissible Gaussian test functions $h^{\mathrm{even}}_{p,\varepsilon}$.
The diagonal of such a matrix would be connected to the
curvature sum $B$ via the identity $O''(\tfrac12)=-2B$, while
its off-diagonal structure encodes the prime-zero coupling studied
in Papers~1--5 of this series.

The Hilbert space $H_W$ derived from the Weil hermitian form
(see Suzuki~\cite{Suzuki25}) is the natural infinite-dimensional
ambient space in which the finite-dimensional matrix
$W_{\kappa,\varepsilon}$ lives as a finite-rank object.
Connes and Consani~\cite{ConnesConsani23} construct the
matrix $\sigma(n,m) = Q_W(\eta_n,\eta_m)$ of the Weil quadratic
form on a basis $\{\eta_n\}$ indexed by integers, using
prolate-spheroidal cutoffs on a compact interval $[1/\lambda,\lambda]$.
A possible future construction of the Weil
transfer operator $W_{\kappa,\varepsilon}$
is indexed by primes
$p\le\kappa$ rather than integers~$n$, uses Gaussian smoothing
of width~$\varepsilon$ on the spectral side rather than a
prolate-spheroidal cutoff, and is connected to the curvature
identity $O''(\tfrac12)=-2B$ established here via its diagonal.
The two constructions are methodologically related but technically
distinct.
\end{remark}

% -------------------------------------------------------
\section*{Non-Circularity}
\label{sec:nocirc}

We document explicitly that the results of this paper do not rely on
the objects they are directed toward.

\textbf{No RH as input.}
At no point is the Riemann Hypothesis assumed.
The ordinates $\gamma_k$ are used as numerically specified inputs, and
no property of their distribution requiring the Riemann Hypothesis is
assumed or used.

\textbf{No GUE, Montgomery, or Hilbert--P\'olya hypothesis.}
{\sloppy
The operator $\Tt$ and its $\sigma$-parametrised extension
$\Tt(\sigma)$ are constructed explicitly from Gaussian-smoothed prime
phasors in the previous papers of this series~\cite{Paper3,Paper4,Paper5}.
No random-matrix hypothesis, no pair-correlation conjecture, and no
Hilbert--P\'olya postulate enters the construction or the proofs.
\par}

\textbf{$\Tt$ is not a Hilbert--P\'olya operator.}
As observed numerically in~\cite{Paper5} on the tested finite grids,
the eigenvalues follow an inverse-ordinate pattern
$\mu_j\approx\CT/\gamma_{k(j)}$, which is incompatible with a
Hilbert--P\'olya interpretation on those numerical diagnostics.
The paper does not attempt a Hilbert--P\'olya realisation.

\textbf{The Guinand--Weil formula is classical.}
Proposition~\ref{prop:gw} is a classical identity valid for every even,
sufficiently rapidly decaying test function; its derivation does not
require the truth of the Riemann Hypothesis.
Theorems~\ref{thm:main-neg} and~\ref{thm:other}, together with the
algebraic identity of Proposition~\ref{prop:curv-bias}, are proved
unconditionally by direct computation.
Theorem~\ref{thm:trunc} verifies numerically that the
$N=100$ truncation tail is negligible at the stated reference parameters.

\textbf{Numerical results are clearly marked.}
Theorems~\ref{thm:gamma},~\ref{thm:ratio-signcross},
and~\ref{thm:intbias}, and Corollary~\ref{cor:curvature}, rely on
numerical evaluation of finite sums or one-dimensional numerical
quadrature; where zero-ordinate truncation enters, the tail is
controlled by Theorem~\ref{thm:trunc}.
Each such statement is labelled [numerical] in its heading.

\textbf{Conditional statements are explicitly flagged.}
The conditional remark following Corollary~\ref{cor:curvature} is
stated conditionally on the stronger exact stationarity condition
$O'(\tfrac12)=0$, which is distinguished from the weaker
first-derivative cancellation target in Problem~\ref{op:stat},
and adds no further unproved input beyond this explicit assumption.
Problem~\ref{op:transfer} is named openly and is not used as a
hypothesis anywhere in \S\S\,\ref{sec:main}--\ref{sec:curvcor}.

\textbf{$\gamma_k$ as inputs, not conclusions.}
The ordinates appear as numerically specified inputs to the model.
The claim $\mathrm{Re}(\rho_k)=\tfrac12$ is not used anywhere; the
reference point $\sigma_0=\tfrac12$ is chosen as the object of study,
not assumed to be the unique zero location.

% -------------------------------------------------------
\section*{Acknowledgments}

The author thanks colleagues and external readers for iterative
critical review across draft versions, whose feedback shaped the
final formulation of the results in this paper.

% -------------------------------------------------------
\section*{Call for Feedback}

This paper is part of an open research initiative on prime--zero
coupling and explicit-formula operators.
All papers in the series are freely available on Zenodo under CC~BY~4.0,
and the accompanying verification code is hosted on GitHub.

The author welcomes critical feedback, corrections, and collaboration:

{\raggedright
Zenodo:\\
\url{https://zenodo.org/search?q=metadata.creators.person_or_org.name\%3A\%22Tehrani\%2C\%20Ulrich\%22}\\
GitHub:\\
\url{https://github.com/utehrani/analysislab-nt}
\par}

Topics of particular interest include: analytical approaches to the
asymptotic pointwise bias of Problem~\ref{op:bias-asymp}; the analytic
stationarity bound of Problem~\ref{op:stat}; uniformity of the
curvature inequality across a larger parameter range
(Problem~\ref{op:uniformity}); the integrated-to-direct transfer of
Problem~\ref{op:transfer}; and, ultimately, the Lagarias--Weil bridge
of Problem~\ref{op:weil-bridge}.

We particularly welcome expert comment on the following question:
for a Gaussian-smoothed family $g_\varepsilon$ with
$g_\varepsilon\to\delta$ as $\varepsilon\to 0$, is it known whether
the sign of $W(g_\varepsilon * g_\varepsilon^\vee)$ at
finite~$\varepsilon$ carries arithmetic information, or is positivity
only meaningful in the distributional limit?

% -------------------------------------------------------
\newpage
\begin{thebibliography}{99}

\bibitem{Paper1}
U.~Tehrani,
``A Curvature Decomposition of the Explicit Formula
for the Riemann Zeta Function'',
Zenodo, \url{https://doi.org/10.5281/zenodo.19025598}, 2026.

\bibitem{Paper2}
U.~Tehrani,
``From Local Curvature to the Weil Functional:
An Explicit Construction'',
Zenodo, \url{https://doi.org/10.5281/zenodo.19106992}, 2026.

\bibitem{Paper3}
U.~Tehrani,
``A Finite-Cutoff Hilbert-Space Model
for Prime--Zero Energy Structure'',
Zenodo, \url{https://doi.org/10.5281/zenodo.19307989}, 2026.

\bibitem{Paper4}
U.~Tehrani,
``A Dual Operator for Prime--Zero Coupling
and a Conditional Proof of Energy Asymmetry'',
Zenodo, \url{https://doi.org/10.5281/zenodo.19364703}, 2026.

\bibitem{Paper5}
U.~Tehrani,
``Spectral Trace Formula and Smoothed Zero Sums:
A Prime--Zero Duality Framework'',
Zenodo, \url{https://doi.org/10.5281/zenodo.19508547}, 2026.

\bibitem{Code}
U.~Tehrani,
\emph{Verification scripts for the series},
GitHub, \url{https://github.com/utehrani/analysislab-nt}.

\bibitem{Lagarias}
J.C.~Lagarias,
``On a positivity property of the Riemann $\xi$-function'',
\emph{Acta Arithmetica} \textbf{89} (1999), 217--234.

\bibitem{Guinand}
A.P.~Guinand,
``A summation formula in the theory of prime numbers'',
\emph{Proc.\ London Math.\ Soc.}\ (2) \textbf{50} (1948), 107--119.

\bibitem{Weil}
A.~Weil,
``Sur les `formules explicites' de la th\'eorie des nombres
premiers'',
\emph{Comm.\ S\'em.\ Math.\ Univ.\ Lund, Suppl.}\ (1952), 252--265.

\bibitem{Hadamard}
J.~Hadamard,
``Sur la distribution des z\'eros de la fonction $\zeta(s)$
et ses cons\'equences arithm\'etiques'',
\emph{Bull.\ Soc.\ Math.\ France} \textbf{24} (1896), 199--220.

\bibitem{Edwards}
H.M.~Edwards,
\emph{Riemann's Zeta Function},
Academic Press, New York, 1974.

\bibitem{Connes}
A.~Connes,
``Trace Formula in Noncommutative Geometry and the Zeros of the
Riemann Zeta Function'',
\emph{Selecta Math.\ (N.S.)} \textbf{5}.1 (1999), 29--106.

\bibitem{BerryKeating}
M.V.~Berry and J.P.~Keating,
``The Riemann Zeros and Eigenvalue Asymptotics'',
\emph{SIAM Review} \textbf{41}.2 (1999), 236--266.

\bibitem{deBranges}
L.~de~Branges,
``The Riemann Hypothesis for Hilbert Spaces of Entire Functions'',
\emph{Bull.\ Amer.\ Math.\ Soc.\ (N.S.)} \textbf{15}.1 (1986),
1--17.

\bibitem{Burnol}
J.-F.~Burnol,
``On Fourier and Zeta(s)'',
\emph{Forum Math.} \textbf{16}.6 (2004), 789--840.

\bibitem{Bombieri}
E.~Bombieri,
``Remarks on Weil's Quadratic Functional in the Theory of Prime
Numbers, I'',
\emph{Rend.\ Lincei Mat.\ Appl.} \textbf{11}.3 (2000), 183--233.

\bibitem{Li}
X.-J.~Li,
``The Positivity of a Sequence of Numbers and the Riemann
Hypothesis'',
\emph{J.\ Number Theory} \textbf{65}.2 (1997), 325--333.

\bibitem{BombieriLagarias}
E.~Bombieri and J.C.~Lagarias,
``Complements to Li's Criterion for the Riemann Hypothesis'',
\emph{J.\ Number Theory} \textbf{77}.2 (1999), 274--287.

\bibitem{RubinsteinSarnak}
M.~Rubinstein and P.~Sarnak,
``Chebyshev's Bias'',
\emph{Experimental Math.} \textbf{3}.3 (1994), 173--197.


\bibitem{IwaniecKowalski04}
H.~Iwaniec and E.~Kowalski,
\emph{Analytic Number Theory},
American Mathematical Society Colloquium Publications,
vol.~53, AMS, Providence, RI, 2004.

\bibitem{ConnesConsani23}
A.~Connes and C.~Consani,
``Spectral triples and zeta-cycles,''
\textit{Enseign.\ Math.} \textbf{69} (2023), no.~1--2, 93--148.
arXiv:2106.01715.
DOI: 10.4171/LEM/1049.

\bibitem{CarneiroLittmannVaaler13}
E.~Carneiro, F.~Littmann, and J.~D.~Vaaler,
``Gaussian subordination for the Beurling--Selberg extremal problem,''
\textit{Trans.\ Amer.\ Math.\ Soc.} \textbf{365} (2013),
3493--3534.
arXiv:1008.4969.
DOI: 10.1090/S0002-9947-2012-05733-X.

\bibitem{CarneiroMilinovichSoundararajan19}
E.~Carneiro, M.~B.~Milinovich, and K.~Soundararajan,
``Fourier optimization and prime gaps,''
\textit{Comment.\ Math.\ Helv.} \textbf{94} (2019), no.~3, 533--568.
arXiv:1708.04122.
DOI: 10.4171/CMH/467.

\bibitem{BalanzarioCardenas25}
E.~P.~Balanzario and D.~E.~C\'{a}rdenas Romero,
``An explicit formula for the zeros of the Riemann zeta function
and the statistics of their local distribution,''
\textit{Ramanujan J.} \textbf{69} (2026), article~21.
DOI: 10.1007/s11139-025-01297-y.

\bibitem{Ramare12}
O.~Ramar\'{e},
``On the behaviour of $\gamma\log p$ modulo~1,''
arXiv:1203.1759 (2012, revised December 2022).
Unpublished preprint.

\bibitem{Bellotti24a}
C.~Bellotti,
``Explicit bounds for the Riemann zeta function and a new
zero-free region,''
\textit{J.\ Math.\ Anal.\ Appl.} \textbf{536} (2024),
no.~2, art.~128249.
arXiv:2306.10680.
DOI: 10.1016/j.jmaa.2024.128249.

\bibitem{Bellotti24b}
C.~Bellotti,
``An explicit log-free zero density estimate for the Riemann
zeta-function,''
\textit{J.\ Number Theory} \textbf{269} (2025), 37--77.
arXiv:2405.12545.
DOI: 10.1016/j.jnt.2024.10.001.

\bibitem{MossinghoffTrudgianYang24}
M.~J.~Mossinghoff, T.~S.~Trudgian, and A.~Yang,
``Explicit zero-free regions for the Riemann zeta-function,''
\textit{Res.\ Number Theory} \textbf{10} (2024), art.~11.
DOI: 10.1007/s40993-023-00498-y.

\bibitem{Suzuki25}
M.~Suzuki,
``On the Hilbert space derived from the Weil distribution,''
\textit{Canad.\ J.\ Math.}, First View (2025), 1--24.
arXiv:2301.00421.
DOI: 10.4153/S0008414X25101739.

\bibitem{BellottiWong24}
C.~Bellotti and P.-J.~Wong,
``Improved estimates for the argument and zero-counting function
of the Riemann zeta-function,''
arXiv:2412.15470v2 (2025).

\bibitem{BrentPlattTrudgian21}
R.~P.~Brent, D.~J.~Platt, and T.~S.~Trudgian,
``Accurate estimation of sums over zeros of the Riemann
zeta-function,''
\textit{Math.\ Comp.} \textbf{90} (2021), no.~332, 2923--2935.
DOI: 10.1090/mcom/3652.
\end{thebibliography}

% -------------------------------------------------------
\enlargethispage{3\baselineskip}
\vspace*{\fill}
\begin{minipage}{\linewidth}
\rule{\linewidth}{0.4pt}\smallskip\\
\small\emph{Paper~6 v1.25 \textperiodcentered\ August~2026
\textperiodcentered\ Pauca sed matura}
\end{minipage}

\end{document}
